\chapter{Opis założeń projektu}
\section{Cel projektu}

Realizacja tego projektu obejmuje następujące cele szczegółowe:
\begin{itemize}
    \item Zapewnienie klientom płynnego i szybkiego procesu odnalezienia interesującej ich pozycji z menu, jej opcjonalnej konfiguracji autorskiej w interaktywnym kreatorze oraz bezproblemowej finalizacji zakupu poprzez moduł koszyka.
    \item Wykreowanie bezpiecznej i czytelnej strefy osobistej (portal ,,Moje konto''), ułatwiającej klientom podgląd historii zrealizowanych zamówień, weryfikację statusu zwrotów oraz zarządzanie kluczowymi danymi profilowymi.
    \item Zbudowanie intuicyjnego pulpitu udostępniającego panel administracyjny (back-office) dla pracowników pizzerii, co znacząco upraszcza zarządzanie asortymentem oraz modyfikowanie karty dań widocznej na froncie sklepu.
\end{itemize}

\section{Zakres funkcjonalności}
System „Wito’s” został zaprojektowany jako kompleksowe rozwiązanie wspierające pełny cykl obsługi zamówienia. Poniżej przedstawiono szczegółowy wykaz funkcjonalności z podziałem na moduły i role użytkowników:
\vspace{10pt}

\textbf{Zarządzanie produktami i koszykiem:}
\begin{itemize}
\item Przeglądanie interaktywnego menu.
\item Dynamiczne dodawanie i usuwanie produktów z koszyka zamówienia.
\item Przegląd podsumowania zamówienia wraz z naliczonymi kosztami dostawy.
\end{itemize}
    
\textbf{Zaawansowany kreator pizzy:}
\begin{itemize}
\item Personalizacja składników według indywidualnych preferencji użytkownika.
\item Funkcja podziału pizzy na dwie różne połowy z niezależnym doborem
dodatków.
\item Możliwość zwiększenia ilości (gramatury) wybranych składników.
\item System walidacji: nakładanie limitów na liczbę składników w celu zachowania standardów technologicznych.
\end{itemize}
    
\textbf{Finalizacja i monitoring:}
\begin{itemize}
\item Obsługa kodów rabatowych.
\item Śledzenie statusu realizacji zamówienia w czasie rzeczywistym.
\end{itemize}

\textbf{Profil użytkownika:}
\begin{itemize}
\item Dostęp do historii zamówień.
\item Zarządzanie danymi adresowymi.
\item Zintegrowany moduł reklamacji i zwrotów.
\end{itemize}

\textbf{Zarządzanie Produkcją (Widok Kuchni):}
\begin{itemize}
\item Podgląd kolejki aktywnych zamówień wraz ze szczegółową specyfikacją składników.
\item Zarządzanie procesem przygotowania poprzez zmianę statusów (np. „w przygotowaniu”, „gotowe”).
\end{itemize}
    
\textbf{Zarządzanie Logistyką (Widok Dostawcy):}
\begin{itemize}
\item Dostęp do listy zamówień gotowych do odbioru wraz z danymi klienta.
\item Aktualizacja statusów doręczenia przesyłki.
\end{itemize}
    
\textbf{Zarządzanie Systemem (Panel Administratora):}
\begin{itemize}
\item Pełny nadzór nad obiegiem zamówień i rozpatrywanie wniosków o zwrot środków.
\item Zarządzanie bazą użytkowników oraz uprawnieniami personelu.
\item Konfiguracja asortymentu: edycja menu oraz aktualizacja cen.
\end{itemize}


